\section{Skills}

\subsection{Identifying Needs}

Needs are sometimes not easy to perceive directly; they are often hidden
behind feelings, judgments, or strategies \citep{rizzolatti2004}.
Here is a systematic three-step method to help you move from feeling to
need:

\begin{enumerate}
    \item \textbf{Step 1: Feel the present moment.}
    When a strong emotion arises (anger, sadness, fear), do not rush to
    act or express.
    Take a deep breath and bring your attention to your body:
    \textbf{``What am I feeling right now?''} (returning to the skills in
    Chapter 2) \citep{decety2021}.

    \item \textbf{Step 2: Ask from feeling to need.}
    Ask yourself: \textbf{``What unmet need is this feeling telling me
    about?''}
    Refer to the Universal Needs Table (Table \ref{tab:universal_needs})
    and try to name the need that is closest to your heart \citep{warneken2018}.

    \item \textbf{Step 3: Distinguish need from strategy.}
    Check whether the ``need'' you found is truly a need or a strategy.
    Ask yourself: \textbf{``Is this a universal desire that anyone would
    have, or a specific thing I want the other person to do?''}
    If it is the latter, keep digging deeper until you touch the real
    need \citep{streubel2021}.
\end{enumerate}

\begin{table}[h]
\centering
\caption{Correspondence between Feelings and Underlying Needs}
\label{tab:feeling_to_need}
\begin{tabular}{|l|p{5cm}|p{7cm}|}
\hline
\textbf{Feeling} & \textbf{Possible Underlying Need} & \textbf{Example} \\ \hline
Loneliness & Connection, companionship, belonging & I long to feel authentically connected \\ \hline
Anger & Fairness, respect, autonomy & I long for my boundaries to be recognized and respected \\ \hline
Fear & Security, certainty, protection & I long to know that I am safe in this relationship \\ \hline
Sadness & Mourning, companionship & I long for my loss to be authentically seen and accompanied \\ \hline
Anxiety & Order, certainty & I long for a clear understanding of what is going to happen \\ \hline
Grievance & Fairness & I long for my efforts and intentions to be correctly understood \\ \hline
Emptiness & Meaning, connection, sense of value & I long for my existence and actions to be meaningful \\ \hline
Jealousy & Belonging, security, self-worth & I long to be cherished in this relationship \\ \hline
Loss & Contribution seen, connection, expectation & I long for my efforts to be felt and recognized \\ \hline
\end{tabular}
\end{table}

\subsection{``What Do You Need?''}

Similar to ``How are you feeling right now?'' from the previous chapter,
cultivating the habit of asking \textbf{``What do you need right now?''}
from a feeling in interpersonal interactions is one of the most powerful
tools in empathetic communication \citep{lobue2022}.

The power of this question lies in:
\begin{itemize}
    \item It guides the other person's attention from ``whose fault is it''
    to ``what do I truly desire in my heart,'' bypassing the battlefield of
    judgment.
    \item It gives the other person the experience of being valued and seen,
    directly triggering oxytocin secretion at the brain level \citep{christou2025}.
    \item It creates a foundation for both parties to find creative
    solutions---because needs can resonate, while strategies can be
    negotiated \citep{ruba2020}.
\end{itemize}

However, this question is only effective after the other person's emotions
have stabilized (after cortisol levels decrease).
As discussed in Chapter 3, in a state of high cortisol, people cannot
engage in deep introspection---they can only see the other person's
``mistake,'' not anyone's needs \citep{aghaziarati2023}.

\paragraph{When the Other Person Cannot Identify Their Own Needs}
Many people have never been taught to recognize their own needs.
They were trained from childhood to ``obey'' or ``suppress,'' never having
the opportunity to ask themselves ``What do I need?''.
When you ask such a person ``What do you need most?'', they may:
\begin{itemize}
    \item Remain silent or bewildered---indicating that they have never
    been allowed or learned to explore their own needs.
    \item Give a strategy instead of a need---``I need you to come home
    early'' (strategy), rather than ``I need companionship'' (need).
    \item Answer with a judgment---``I need you to stop being so selfish''
    (judgment), rather than ``I need to be valued'' (need).
\end{itemize}

At this time, you can gently assist the other person in translating:
\begin{itemize}
    \item ``When you say you need me to come home early, I think I hear
    that you need more companionship?''
    \item ``When you say you need me to stop being so selfish, I think I
    hear that you need to feel valued and considered?''
\end{itemize}

This \textbf{``Translational Empathy''} is a high-level skill; it is not
to correct the other person, but to help the other person's brain find
deeper language \citep{dunn2000}.

\begin{keyinsight}
Identifying and expressing needs involves systematically moving past
immediate emotional reactions and distinguishing between specific strategies
and fundamental human desires.
When we actively translate judgments or vague requests into the clear language
of needs, we defuse conflict and invite profound collaboration and connection.
\end{keyinsight}
